\section*{Appendix B. Symbol Normalization for the Late Readout Layers}
\addcontentsline{toc}{section}{Appendix B. Symbol Normalization for the Late Readout Layers}

This appendix fixes the notation used in the late Companion expansion layers so that
historical version labels and structural symbols remain auditable across the v2.4.0--v2.11.0
chain.

\paragraph{Historical section labels.}
The blocks entitled v2.4.0 through v2.11.0 are to be read as historical expansion
layers. Their section headings record the version in which the corresponding layer was
introduced, not merely the current document version.

\paragraph{Ledger matrix.}
The channel ledger matrix is denoted by
\[
\Lambda_{\mathrm{ch}}(\Psi)
=
\begin{pmatrix}
\ell_{\mathrm{tri}} & \delta_{\mathrm{tri}} & r_{\mathrm{tri}}\\
\ell_{\perp} & \delta_{\perp} & r_{\perp}\\
\ell_{\mathrm{meas}} & \delta_{\mathrm{meas}} & r_{\mathrm{meas}}\\
\ell_{\mathrm{res}} & \delta_{\mathrm{res}} & r_{\mathrm{res}}
\end{pmatrix}.
\]
It is a bookkeeping object for channel-resolved tangential, normal, and compensator-facing
sizes.

\paragraph{Lyapunov ledger functional.}
The weighted scalar functional introduced in v2.6.0 keeps the notation
\(\mathcal L\). It is a Lyapunov-type control quantity and is not to be confused with the
matrix-valued ledger \(\Lambda_{\mathrm{ch}}\).

\paragraph{Budget totals.}
The aggregated scalar budgets remain
\(\mathfrak B\) for total visible load, \(\mathfrak N\) for normal leakage, and
\(\mathfrak R\) for compensator-facing storage.

\paragraph{Readout and fusion objects.}
The readable-mode trace mass, reconstruction operators, observable extraction maps,
and fusion-defect quantities introduced in v2.9.0--v2.11.0 are interpreted as
structure-preserving descendants of the ledger/budget layer and therefore inherit the
normalization fixed here.



\subsection*{Session-budget rule}
A further practical restriction is useful at the present stage: each local mathematics-first pass should also observe a bounded session budget. The purpose of this rule is not to slow work down, but to prevent a single pass from silently expanding from a targeted closure attempt into a diffuse rewrite wave. In concrete terms, a local subversion should normally do only one of the following as its primary task: (i) stabilize definitions for one A-priority block, (ii) discharge one dependency layer for that block, (iii) perform one focused reference/notation repair required by that block, or (iv) close one clearly delimited appendix-tool support step. If more than one such task would be required, the safer action is to end the current pass, classify the result, and continue by explicit next local subversion.

The session-budget rule therefore complements the anti-parallelization and re-entry rules. A pass is considered well-scoped if its main closure claim can be summarized in one sentence and audited against one active block. A pass should be considered over budget if it begins to change multiple theorem-status boundaries, broad notation regimes, or several major sections at once. In that case, the correct response is not to continue accumulating edits, but to stop, audit, and reopen the next step as a fresh local subversion. This rule is intentionally conservative: its goal is to make every local version legible as a finite mathematical move rather than as an opaque mass of partially mixed intentions.



\subsection*{Archived final release-decision template for v2.27.0}
At this stage the document no longer needs another diffuse local steering note, but a compact decision template that makes explicit under which reading the next public release is justified. The present template distinguishes three cases. In the first, the document is strong enough that the next action is \emph{release now}: this applies if the A-priority blocks are structurally frozen, the gate conditions remain green, and the remaining work consists only of optional proof-density enrichment. In the second, the document should \emph{hold locally for proof-density enrichment}: this applies if the structural spine is already in place, but one deliberately wants a denser local unfolding before exposing the next public release. In the third, the document must \emph{hold locally for unresolved blocker}: this applies only if a genuine release-critical dependency, hypothesis-control issue, or reference-instability problem is still active.

Under the current reading, the first A-block no longer counts as such a blocker. Its local dependency spine has been extended through admissible transport, matrix-sector realization, and admissible trace/branch persistence, and it has already been reconciled with the mathematics-first gate logic. Hence the residual caution is no longer structural incompleteness, but only whether one prefers to accept the present medium proof density or to invest one further local pass in additional unfolding. Accordingly, the document is now close to the decision boundary where a direct step to v2.27.0 becomes reasonable.

\paragraph{Decision states.}
\begin{itemize}[leftmargin=2em]
  \item \textbf{Release now:} the first A-block is structurally release-supporting, no primary blocker remains, and only optional densification is deferred.
  \item \textbf{Hold locally for proof-density enrichment:} no structural blocker remains, but one more local pass is preferred for denser exposition.
  \item \textbf{Hold locally for unresolved blocker:} a genuine reference, dependency, or hypothesis-control failure is still present.
\end{itemize}

\paragraph{Current reading.}
At the present local stage, the third state no longer appears to apply to the first A-block. The second and first states are now the relevant ones: either one accepts the current medium-density closure as sufficient for the next public release, or one performs at most one more local enrichment pass before the public step.

\paragraph{Release consequence.}
The practical consequence is that any further local subversion should no longer open a new mathematical front. It should either confirm release-readiness or add one last density-improving clarification and then stop.




\subsection*{2J. Open-questions closure pass and chain check}
The present stage is also strong enough to classify several previously open questions more explicitly. The purpose of this pass is not to pretend that every weak-density zone has disappeared, but to distinguish between (i) questions that can already be closed reliably at the current structural level, (ii) questions that remain structurally intact but only under an explicit assumption or still-thin proof layer, and (iii) questions that should remain openly deferred. This distinction is especially important in a structure-first setting, because a locally thin argument is not automatically broken: some passages remain admissible as conditional or contradiction-style steps, provided that their status is made explicit and no illegitimate theorem-promotion occurs.

\paragraph{Closure classification.}
At the current stage the following questions may be treated as \emph{closable now}: the dependency spine from ground condition to HC and confinement; the compatibility of closure with projector stabilization; the admissible transport continuation on the stabilized sector; the admissible-sector action of the Millennium matrix; and the persistence of admissible trace/branch continuation once the stabilized sector has been fixed. These are not yet maximally unfolded everywhere, but they are no longer merely heuristic. By contrast, the following questions are presently better read as \emph{intact but proof-density limited}: the strongest available mathematical reading of the Millennium matrix as converter/intertwiner; the full density of trace/branch continuation lemmas; and the extent to which some still-compact structural passages should already be promoted beyond proposition level. Finally, the following remain \emph{intentionally open}: any full universality claim for the computational strand, any finished proof-kernel or compiler reading, and any interpretation that would make the mathematics depend on the computational teaser architecture.

\paragraph{Chain-check statuses.}
To keep the proof chain honest without overcalling local weakness as failure, the document now adopts four chain-check statuses for future mathematical passes:
\begin{itemize}[leftmargin=2em]
  \item \textbf{intact:} the dependency chain carries without visible status or continuity defect;
  \item \textbf{intact, but proof-density weak:} the chain carries, but one or more transitions still need denser local unfolding;
  \item \textbf{intact under explicit assumption:} the argument is acceptable only because an assumption, contradiction setup, or provisional hypothesis is openly declared;
  \item \textbf{blocker:} a genuine dependency break, status violation, or contradiction prevents closure.
\end{itemize}

\paragraph{Current reading.}
Under this classification, the first major dependency spine is currently best read as \emph{intact, but proof-density weak} rather than as blocked. No primary structural contradiction is visible in the path from ground condition through HC, closure, admissible transport, Millennium-matrix realization, and admissible trace/branch persistence. The remaining caution is therefore mainly one of density and explicitness, not one of broken continuity.

\paragraph{Local consequence.}
The practical consequence is conservative but positive. Several earlier open questions may already be marked as \emph{reliably closable at the present stage}, provided that the document does not silently promote weak-density passages to stronger theorem status than the current support allows. The next mathematics-first passes should therefore use the chain-check statuses as a standing audit device: strengthen density where useful, keep assumption-based passages openly marked, and treat only genuine continuity failure as a blocker.




\subsection*{2K. Tool extraction and leverage-point scan}
The next mathematics-first pass should not treat every locally thin argument as something that must immediately be expanded line by line. In a structure-first document, some still-wobbly passages already contain reusable mathematical moves that can be extracted earlier as \emph{tools}, provided that their status is made explicit. The present scan therefore distinguishes between (i) tools that are already operationally useful now, even if their surrounding proof density is still moderate, (ii) tools that remain usable only under an explicit structural assumption, and (iii) a small number of leverage points whose clarification would simplify several later proof passages at once.

\paragraph{Provisional tool extraction.}
At the current stage the following reusable tools may already be extracted from the first dependency spine. First, a \emph{closure-transfer tool}: once the admissible sector has been fixed, closure may be transported across projector stabilization without reopening the sector boundary. Second, a \emph{projected transport tool}: admissible transport may be read modulo the stabilized projector rather than by re-proving closure from scratch after each transport step. Third, a \emph{matrix-sector tool}: the Millennium matrix may already be used as an admissible-sector-preserving operator on the stabilized transported sector. Fourth, a \emph{trace/branch persistence tool}: admissible continuation in the trace/branch layer may be used without reintroducing an external admissibility source. None of these tools should yet be oversold as maximal-density theorems everywhere; however, they are already strong enough to simplify downstream arguments by replacing repeated local closure checks with a smaller number of reusable invariance moves.

\paragraph{Status discipline for extracted tools.}
To keep the extraction mathematically honest, each such tool should be marked in one of three ways:
\begin{itemize}[leftmargin=2em]
  \item \textbf{operational tool:} reusable now, with the present support level;
  \item \textbf{assumption-backed tool:} reusable only under an explicitly declared hypothesis or contradiction setup;
  \item \textbf{blocked tool candidate:} not yet safe to reuse because a real dependency gap remains.
\end{itemize}
Under the current chain check, the closure-transfer, projected transport, matrix-sector, and trace/branch persistence tools are best read as \emph{operational tools}, though some still sit in the ``proof-density weak'' zone rather than the fully unfolded zone.

\paragraph{Leverage-point scan.}
The present scan suggests that the most useful immediate leverage point is not a new broad theorem family, but a single stronger organizing hypothesis: namely that the admissible stabilized sector behaves as the true fixed working domain for closure, transport, matrix action, and admissible continuation. Even if one initially formulates this only as a hypothesis-backed simplifier, it has unusually high leverage: if strengthened, it compresses several later arguments by replacing repeated local admissibility restoration with one global working-domain principle. In other words, many later proofs become shorter once they are allowed to start with ``assume the admissible stabilized sector has already been fixed'' and then work internally on that domain.

\paragraph{Use of the computational layer as a secondary aid.}
The computational chapter may already help here, but only in a subordinate role. MML/MMA/MCP should not drive the mathematics; however, the language layer can already serve as a \emph{notation and workflow aid} for tool extraction. In particular, compact MML-style descriptions can help identify whether a candidate move is really one reusable operation---for example \texttt{project -> transport -> trace}---or whether it still hides an unresolved mathematical gap. Thus the computational layer may already be used as a diagnostic simplifier for tool extraction, while the mathematical status of the extracted tool remains decided purely on the mathematics-first side.

\paragraph{Current reading and local consequence.}
The current reading is therefore positive. The first dependency spine is no longer merely a place where one asks whether the chain survives; it is already rich enough to yield a first family of reusable structural tools. At the same time, the highest-leverage next target is also visible: stabilize the admissible working-domain principle as far as possible, because this one point appears to simplify closure, transport, matrix action, and trace/branch continuation simultaneously. Any next local pass should therefore prefer a leverage-point clarification over broad diffuse proof growth.



